
\usepackage[xetex,top=24mm,bottom=25mm,left=27mm,right=25mm,headheight=15pt]{geometry}
\usepackage{fontspec,xeCJK}
\setmainfont{Times New Roman}
\setCJKmainfont[BoldFont=SimHei,ItalicFont=KaiTi]{SimSun}
\setCJKsansfont{Microsoft YaHei}
\setCJKmonofont{SimSun}
\usepackage{amsmath,amssymb,amsthm,mathtools}
\usepackage{booktabs,longtable,array,graphicx}
\usepackage{fancyhdr}
\usepackage{xcolor}
\usepackage[unicode,bookmarks=false,colorlinks=true,linkcolor=blue!45!black,citecolor=blue!45!black,urlcolor=blue!45!black]{hyperref}
\newcommand{\Z}{\mathbb Z}
\newcommand{\Q}{\mathbb Q}
\newcommand{\F}{\mathbb F}
\newcommand{\R}{\mathbb R}
\newcommand{\SL}{\operatorname{SL}}
\newcommand{\GL}{\operatorname{GL}}
\newcommand{\diag}{\operatorname{diag}}
\newcommand{\rad}{\operatorname{rad}}
\newcommand{\cont}{\operatorname{cont}}
\newcommand{\Um}{\operatorname{Um}}
\newcommand{\ord}{\operatorname{ord}}
\newtheorem{theorem}{定理}[chapter]
\newtheorem{lemma}[theorem]{引理}
\newtheorem{proposition}[theorem]{命题}
\newtheorem{corollary}[theorem]{推论}
\theoremstyle{definition}
\newtheorem{definition}[theorem]{定义}
\newtheorem{example}[theorem]{例}
\newtheorem{remark}[theorem]{说明}
\renewcommand{\proofname}{证明}
\renewcommand{\contentsname}{目录}
\renewcommand{\bibname}{参考文献}
\renewcommand{\tablename}{表}
\renewcommand{\figurename}{图}
\renewcommand{\chaptername}{讲}
\renewcommand{\appendixname}{附录}
\numberwithin{equation}{chapter}
\setcounter{tocdepth}{1}
\setcounter{secnumdepth}{2}
\setlength{\parindent}{2em}
\setlength{\parskip}{4pt plus 1pt}
\linespread{1.16}
\emergencystretch=3em
\allowdisplaybreaks[2]
\raggedbottom
\pagestyle{fancy}
\fancyhf{}
\fancyhead[L]{\small 素数平均：从例子到完整证明}
\fancyhead[R]{\small\thepage}
\renewcommand{\headrulewidth}{0.3pt}
\fancypagestyle{plain}{\fancyhf{}\fancyfoot[C]{\thepage}\renewcommand{\headrulewidth}{0pt}}
\makeatletter
\def\@makechapterhead#1{\vspace*{8mm}{\parindent\z@\raggedright\normalfont
 \ifnum\c@secnumdepth>\m@ne\if@mainmatter
 {\Large\sffamily\color{blue!40!black}\@chapapp\ \thechapter\par}\vskip12pt\fi\fi
 \interlinepenalty\@M\LARGE\bfseries #1\par\nobreak\vskip24pt}}
\def\@makeschapterhead#1{\vspace*{8mm}{\parindent\z@\raggedright\normalfont
 \interlinepenalty\@M\LARGE\bfseries #1\par\nobreak\vskip24pt}}
\makeatother
\hypersetup{pdftitle={素数平均：从例子到完整证明},pdfsubject={中文数学讲义}}
